\section*{अध्याय \ref{ch:g11:enzymes-and-phenotype} --- एंजाइम और लक्षणप्ररूप}

\begin{solution}{exo:g11:enzymes-and-phenotype:1}
\omterm{def:g11:enzymes-and-phenotype:phenotype}{जीनप्ररूप}: ढोए हुए विकल्पियों का समुच्चय। \omterm{def:g11:enzymes-and-phenotype:phenotype}{लक्षणप्ररूप}: देखे जा सकने वाले
अभिलक्षणों का समुच्चय, आणविक, कोशिकीय और जीव के पैमाने पर।
\end{solution}

\begin{solution}{exo:g11:enzymes-and-phenotype:2}
\omterm{def:g11:enzymes-and-phenotype:enzyme}{एंजाइम} की वह जेब जहाँ \omterm{def:g11:enzymes-and-phenotype:enzyme}{क्रियाधार} बँधता है और बदला जाता है। विशिष्टता: वह
स्थल एक ही \omterm{def:g11:enzymes-and-phenotype:enzyme}{क्रियाधार} (या \omterm{def:g11:enzymes-and-phenotype:enzyme}{क्रियाधारों} के एक ही कुल) से मेल खाता है और एक ही
अभिक्रिया उत्प्रेरित करता है।
\end{solution}

\begin{solution}{exo:g11:enzymes-and-phenotype:3}
\qty{37}{\celsius} के पास; और \qty{45}{\celsius} पर लगभग एक-तिहाई बचता है।
\end{solution}

\begin{solution}{exo:g11:enzymes-and-phenotype:4}
आणविक: टायरोसिनेज़ अनुपस्थित या निष्क्रिय, कोई मेलानिन नहीं बनता। कोशिकीय:
वर्णक \omterm{def:g10:cells-common-unit:cell}{कोशिकाएँ} मौजूद पर मेलानिन से ख़ाली। जीव का: सफ़ेद बाल, बहुत फीकी
त्वचा और आँखें, तथा प्रकाश और धूप की जलन के प्रति संवेदनशीलता।
\end{solution}

\begin{solution}{exo:g11:enzymes-and-phenotype:5}
पेप्सिन का मोड़ केवल तेज़ अम्ल में (pH 2 पर) काम करता है; और आँत उदासीन से
ज़रा क्षारीय है, जहाँ वक्र शून्य सक्रियता दिखाता है।
\end{solution}

\begin{solution}{exo:g11:enzymes-and-phenotype:6}
एक काम करता \omterm{def:g10:universal-dna:gene}{युग्मविकल्पी} सामान्य से आधा \omterm{def:g11:enzymes-and-phenotype:enzyme}{एंजाइम} बनाता है, पर \omterm{def:g11:enzymes-and-phenotype:enzyme}{एंजाइम} का एक अणु
प्रति सेकंड हज़ारों बार काम करता है: इसलिए आधी मात्रा भी सामान्य वर्णक
बनाने भर टायरोसीन बदल देती है। \omterm{def:g11:enzymes-and-phenotype:phenotype}{लक्षणप्ररूप} संतृप्त है; और केवल दोनों
विकल्पियों के चूकने पर ही वह दिखता है।
\end{solution}

\begin{solution}{exo:g11:enzymes-and-phenotype:7}
उनकी वर्णकहीनता दो अलग-अलग \omterm{def:g10:universal-dna:gene}{जीनों} की रुकावटों से आती है (पथ के दो अलग-अलग
\omterm{def:g11:enzymes-and-phenotype:enzyme}{एंजाइम}, या मेलानिन के परिवहन के)। बच्चे को हर \omterm{def:g10:universal-dna:gene}{जीन} का एक काम करता \omterm{def:g10:universal-dna:gene}{युग्मविकल्पी} उस
जनक से मिलता है जिसकी रुकावट कहीं और है, इसलिए दोनों \omterm{def:g11:enzymes-and-phenotype:enzyme}{एंजाइम} काम करते हैं और
वर्णक बन जाता है।
\end{solution}

\begin{solution}{exo:g11:enzymes-and-phenotype:8}
फेनिलऐलानिन जमा होता है; और टायरोसीन (तथा उससे निकलने वाली चीज़ें) की कमी
रहती है। चूँकि मेलानिन टायरोसीन से बनता है, इसलिए उसकी कमी वर्णक घटा देती
है: बिना इलाज वाले फेनिलकीटोनूरिया में हल्के बाल और त्वचा आम हैं।
\end{solution}

\begin{solution}{exo:g11:enzymes-and-phenotype:9}
जहाँ दूध मुख्य आहार हो, वहाँ उसे पचाना बिना बीमारी के ऊर्जा तथा कैल्शियम
देता है, इसलिए वह बना रहने वाला \omterm{def:g10:universal-dna:gene}{युग्मविकल्पी} एक लाभ है। और जहाँ वयस्क दूध पीते
ही नहीं, वहाँ वह \omterm{def:g10:universal-dna:gene}{युग्मविकल्पी} कुछ नहीं बदलता; वह केवल उसी पर्यावरण में ``अधिक
स्वस्थ'' है जो उसका \omterm{def:g11:enzymes-and-phenotype:enzyme}{क्रियाधार} देता हो।
\end{solution}

\begin{solution}{exo:g11:enzymes-and-phenotype:10}
टायरोसिनेज़ के लिए: आणविक, \omterm{def:g11:enzymes-and-phenotype:enzyme}{क्रियाधार} अब बैठता ही नहीं, कोई मेलानिन नहीं;
कोशिकीय, ख़ाली वर्णक \omterm{def:g10:cells-common-unit:cell}{कोशिकाएँ}; जीव का, वर्णकहीनता। लैक्टेज़ के लिए:
लैक्टोस बिना पचा; आँत की \omterm{def:g10:cells-common-unit:cell}{कोशिकाएँ} उस शर्करा से वंचित और बृहदांत्र उसका
\omterm{def:g10:cell-metabolism:fermentation}{किण्वन} करता हुआ; यानी असहिष्णुता।
\end{solution}

\begin{solution}{exo:g11:enzymes-and-phenotype:11}
उसके सिरे \qty{35}{\celsius} से गरम बने रहते हैं, \omterm{def:g11:enzymes-and-phenotype:enzyme}{एंजाइम} मुश्किल से ही काम
करता है, और वह बच्चा पूरा का पूरा फीका रहता है। जाड़े में बाहर रहने पर कान
और पंजे ठंडे पड़ते हैं, \omterm{def:g11:enzymes-and-phenotype:enzyme}{एंजाइम} सक्रिय हो जाता है, और वहाँ दोबारा उगते रोएँ
गहरे रंग के आते हैं।
\end{solution}

\begin{solution}{exo:g11:enzymes-and-phenotype:12}
आणविक: आधा हीमोग्लोबिन हँसिया रूप वाला है। कोशिकीय: सामान्य ऑक्सीजन पर
\omterm{def:g10:cells-common-unit:cell}{कोशिकाएँ} गोल बनी रहती हैं; और बहुत कम ऑक्सीजन पर कुछ हँसिया बन जाती हैं।
जीव का: रोज़मर्रा के जीवन में स्वस्थ, पर ऊँचाई पर या गोताख़ोरी में ख़तरे
में --- और मलेरिया से बचाव। वाहक का \omterm{def:g11:enzymes-and-phenotype:phenotype}{लक्षणप्ररूप} बोझ है या लाभ, यह
पर्यावरण (ऑक्सीजन, मलेरिया) तय करता है।
\end{solution}

\begin{solution}{exo:g11:enzymes-and-phenotype:13}
\qty{41}{\celsius} पर \omterm{def:g11:enzymes-and-phenotype:enzyme}{एंजाइम} अपनी अधिकांश सक्रियता बनाए रखते हैं (वक्र
अपनी चोटी के पास है); और \qty{43}{\celsius} पर, जहाँ \omterm{def:g11:gene-expression:protein}{प्रोटीन} खुलने लगते
हैं, सक्रियता तेज़ी से गिरती है, और यदि वह खुलना आगे बढ़े तो अपलट्य है:
\omterm{def:g11:enzymes-and-phenotype:enzyme}{एंजाइम} धीमे नहीं पड़ते, नष्ट हो जाते हैं।
\end{solution}

\begin{solution}{exo:g11:enzymes-and-phenotype:14}
प्रति \omterm{def:g10:muscles-and-joints:muscle}{पेशी तंतु} अधिक \omterm{def:g10:cells-common-unit:organelle}{सूत्रकणिकाएँ} और केशिकाएँ; बड़े प्रघात आयतन वाला बड़ा
हृदय; प्रति लीटर रक्त अधिक हीमोग्लोबिन; बड़े \omterm{def:g10:exercise-and-energy:glycogen}{ग्लाइकोजन} भंडार। इनमें से हर
एक वही \omterm{def:g11:enzymes-and-phenotype:phenotype}{जीनप्ररूप} किसी अलग पर्यावरण में देकर बना \omterm{def:g11:enzymes-and-phenotype:phenotype}{लक्षणप्ररूप} है:
``आनुवंशिक'' लक्षण वह परास हैं जिसकी छूट \omterm{def:g11:enzymes-and-phenotype:phenotype}{जीनप्ररूप} देता है, कोई जमा हुआ मान
नहीं।
\end{solution}

\begin{solution}{exo:g11:enzymes-and-phenotype:15}
\omterm{def:g11:enzymes-and-phenotype:phenotype}{जीनप्ररूप} तय करता है कि कौन-सा \omterm{def:g11:enzymes-and-phenotype:enzyme}{एंजाइम} बनेगा और वह कैसा बरतेगा: \omterm{def:g11:enzymes-and-phenotype:enzyme}{एंजाइम} 1
नहीं (फेनिलकीटोनूरिया), टायरोसिनेज़ नहीं (वर्णकहीनता), ताप-संवेदी
टायरोसिनेज़ (स्याम)। पर्यावरण नतीजा केवल उतना ही बदल सकता है जितनी छूट वह
\omterm{def:g11:enzymes-and-phenotype:enzyme}{एंजाइम} देता है: आहार \omterm{def:g11:enzymes-and-phenotype:enzyme}{क्रियाधार} हटाकर मस्तिष्क बचा लेता है, पर ग़ायब \omterm{def:g11:enzymes-and-phenotype:enzyme}{एंजाइम}
बना नहीं सकता; ठंड स्याम वाले \omterm{def:g11:enzymes-and-phenotype:enzyme}{एंजाइम} को चला देती है, पर वर्णकहीन वाले को
कोई भी ताप नहीं चलाता; और पर्यावरण की कोई भी चीज़ वर्णकहीन को उसका वर्णक
नहीं देती। \omterm{def:g11:enzymes-and-phenotype:phenotype}{जीनप्ररूप} परास तय करता है; पर्यावरण उसी के भीतर चुनता है।
\end{solution}

\begin{solution}{pb:g11:enzymes-and-phenotype:1}
\textbf{1.} धड़ \qty{38}{\celsius}: लगभग 2\%, फीका। टाँगें
\qty{36}{\celsius}: लगभग 20\%, फीकी। पंजे \qty{33}{\celsius}: 80\%,
गहरे। कान और थूथन \qty{31}{\celsius}: 90\% से ऊपर, गहरे।

\textbf{2.} फीका शरीर और फीकी टाँगें, गहरे पंजे, कान, थूथन (और पूँछ):
यानी स्याम वाला ढर्रा।

\textbf{3.} हर जगह 100\%: यानी एक-सी वर्णकयुक्त बिल्ली।

\textbf{4.} गर्भाशय में बच्चे का हर हिस्सा \qty{38}{\celsius} पर होता है,
\omterm{def:g11:enzymes-and-phenotype:enzyme}{एंजाइम} हर जगह निष्क्रिय रहता है, और जन्म से पहले कोई वर्णक जमा ही नहीं
होता; और जन्म के बाद सिरे ठंडे पड़ते ही वे बिंदु गहरे होने लगते हैं।

\textbf{5.} लगभग \qty{35.5}{\celsius} पर, जहाँ सक्रियता 30\% को पार करती
है: और उस सरहद को 33 तथा \qty{37}{\celsius} के बीच वक्र का तीखा उतार तय
करता है।

\textbf{6.} गहरा: \qty{30}{\celsius} पर उगते बालों में \omterm{def:g11:enzymes-and-phenotype:enzyme}{एंजाइम} पूरी तरह
सक्रिय है।

\textbf{7.} फीका: \qty{38}{\celsius} पर \omterm{def:g11:enzymes-and-phenotype:enzyme}{एंजाइम} निष्क्रिय है।

\textbf{8.} वर्णक बाल में उसके उगते समय ही जमा होता है और बाद में न हटाया
जा सकता है न जोड़ा; और रंग तभी बदलता है जब वे बाल झड़ जाएँ और स्थानीय ताप
पर उगे नए बाल उनकी जगह लें।

\textbf{9.} उस टुकड़े की \omterm{def:g10:cells-common-unit:cell}{कोशिकाओं} का \omterm{def:g11:enzymes-and-phenotype:phenotype}{जीनप्ररूप} पहले और बाद में वही है, और
धड़ के बाक़ी हिस्से जैसा भी; बदला केवल \omterm{def:g11:enzymes-and-phenotype:enzyme}{एंजाइम} का ताप, और उसके साथ वर्णक।
यानी क्रिया \omterm{def:g11:enzymes-and-phenotype:enzyme}{एंजाइम} की सक्रियता पर हुई, जो एक आणविक घटना है।

\textbf{10.} बिल्ली की सारी \omterm{def:g10:cells-common-unit:cell}{कोशिकाएँ} एक ही अंडाणु से उतरी हैं और वही
\omterm{def:g10:universal-dna:gene}{युग्मविकल्पी} ढोती हैं; और मूँड़े हुए टुकड़े वाले प्रयोग किसी क्षेत्र का रंग उसकी
\omterm{def:g10:cells-common-unit:cell}{कोशिकाएँ} बदले बिना बदल देते हैं --- और किसी \omterm{def:g10:universal-dna:gene}{युग्मविकल्पी} को बर्फ़ की थैली से
बदला नहीं जा सकता।

\textbf{11.} आणविक: ऐसा टायरोसिनेज़ जो केवल लगभग \qty{35}{\celsius} से
नीचे काम करता है। कोशिकीय: सिरों पर मेलानिन से भरी वर्णक \omterm{def:g10:cells-common-unit:cell}{कोशिकाएँ}, धड़ में
ख़ाली। जीव का: गहरे बिंदुओं वाली एक फीकी बिल्ली।

\textbf{12.} वह \omterm{def:g10:chemistry-of-life:families}{ऐमीनो अम्ल} वहाँ बैठा है जहाँ वह उस मोड़ को थामे रखता है;
और बदली हुई जगह वह मोड़ कम ताप पर तो जोड़े रखती है पर कुछ ही अंश ऊपर उसे
ढीला पड़ने देती है, इसलिए \omterm{def:g11:enzymes-and-phenotype:enzyme}{सक्रिय स्थल} साधारण \omterm{def:g11:enzymes-and-phenotype:enzyme}{एंजाइम} से पहले ही अपना आकार खो
देता है।

\textbf{13.} साधारण \omterm{def:g10:universal-dna:gene}{युग्मविकल्पी} का \omterm{def:g11:enzymes-and-phenotype:enzyme}{एंजाइम} शरीर के ताप पर पूरी तरह सक्रिय है;
और सामान्य से आधी मात्रा भी हर जगह वर्णक बनाने से कहीं अधिक है। इसलिए
स्याम \omterm{def:g10:universal-dna:gene}{युग्मविकल्पी} उसके सामने अप्रभावी है।

\textbf{14.} आणविक: किसी भी ताप पर कोई काम करता \omterm{def:g11:enzymes-and-phenotype:enzyme}{एंजाइम} नहीं, जबकि दूसरे
में शर्त के साथ काम करता \omterm{def:g11:enzymes-and-phenotype:enzyme}{एंजाइम} है। कोशिकीय: कहीं भी मेलानिन नहीं, जबकि
दूसरे में ठंडे क्षेत्रों में मेलानिन। जीव का: गुलाबी आँखों वाली पूरी सफ़ेद
बिल्ली, जबकि दूसरी में बिंदुओं वाला ढर्रा। वर्णकहीन के लिए ठंड कुछ नहीं
बदलती।

\textbf{15.} यह कि वे उसी \omterm{def:g10:universal-dna:gene}{जीन} का कोई ताप-संवेदी \omterm{def:g10:universal-dna:gene}{युग्मविकल्पी} ढोते हैं; और शायद
\omterm{def:g11:enzymes-and-phenotype:enzyme}{एंजाइम} को अस्थिर करने वाला कोई वैसा ही प्रतिस्थापन, जो हर \omterm{def:g10:biodiversity-scales:species}{जाति} में
स्वतंत्र रूप से उठा।

\textbf{16.} बीस गुना; और \omterm{def:g11:enzymes-and-phenotype:enzyme}{एंजाइम} 1 ग़ायब है, जो फेनिलऐलानिन को टायरोसीन
में बदलता है।

\textbf{17.} अतिरिक्त फेनिलऐलानिन ख़ास तौर पर विकसित होते मस्तिष्क की
\omterm{def:g10:cells-common-unit:cell}{कोशिकाओं} में दख़ल देता है (दूसरे \omterm{def:g10:chemistry-of-life:families}{ऐमीनो अम्लों} की उनकी आपूर्ति में और उनके
परिपक्व होने में); त्वचा की \omterm{def:g10:cells-common-unit:cell}{कोशिकाएँ} उसे सह लेती हैं। आणविक ख़राबी हर जगह
वही है; कोशिकीय नतीजा नहीं।

\textbf{18.} \num{1200} का पाँचवाँ हिस्सा \qty{240}{\micro mol/L} है, जो
\num{360} से नीचे है: यानी सुरक्षित। \omterm{def:g11:enzymes-and-phenotype:enzyme}{एंजाइम} के बिना रक्त का स्तर बस सेवन
के पीछे-पीछे चलता है, इसलिए \omterm{def:g11:enzymes-and-phenotype:phenotype}{लक्षणप्ररूप} अकेले ग़ायब \omterm{def:g11:enzymes-and-phenotype:enzyme}{एंजाइम} से नहीं, बल्कि
\omterm{def:g11:enzymes-and-phenotype:enzyme}{क्रियाधार} की आपूर्ति --- यानी पर्यावरण --- से तय होता है।

\textbf{19.} शरीर के \omterm{def:g11:gene-expression:protein}{प्रोटीन} बनाने के लिए फेनिलऐलानिन चाहिए ही और उसे बनाया
नहीं जा सकता; बिलकुल न मिले तो वृद्धि रुक जाती है। इसलिए सेवन \omterm{def:g11:gene-expression:protein}{प्रोटीन}
संश्लेषण भर होना चाहिए और उससे अधिक नहीं, और उसके लिए नापना तथा समायोजित
करना पड़ता है।

\textbf{20.} \omterm{def:g11:enzymes-and-phenotype:phenotype}{जीनप्ररूप} ने \omterm{def:g11:enzymes-and-phenotype:enzyme}{एंजाइम} तय किया (बिल्ली में ताप-संवेदी, बच्चे में
अनुपस्थित); पर्यावरण ने उसकी शर्तें तय कीं (त्वचा का ताप; आहार का
फेनिलऐलानिन); और \omterm{def:g11:enzymes-and-phenotype:phenotype}{लक्षणप्ररूप} किसी दिए हुए पर्यावरण में अभिव्यक्त हुए
\omterm{def:g11:enzymes-and-phenotype:phenotype}{जीनप्ररूप} का नतीजा है।
\end{solution}
